\chapter{Background and Literature Review}
\label{ch:literature}

\section{Energy, carbon and the measurement boundary}

\subsection{Energy efficiency and hardware control}
Energy efficiency reduces electricity consumption; carbon efficiency can also improve by using a lower-carbon supply. HPC surveys distinguish scheduling, application optimisation, frequency scaling and power capping, and discuss the difficulty of modelling heterogeneous processors and accelerators \citep{czarnul2019energy,obrien2017survey}. This study changes a fixed workload's admission timing. Device power limits remain unchanged, and lower workload energy is not assumed.

Zeus changes batch size and GPU power limits in recurring DNN jobs~\citep{you2023zeus}. It considers energy and training time because lower power may lengthen execution without reducing total energy. Stanage admission instead retains resources and runtimes, changing energy's timing in the scenario model rather than seeking Zeus's application-level reduction.

Chase adapts GPU power limits to forecast carbon intensity during execution, changing a running job's training progress and energy use~\citep{yang2023chase}. Admission shifting changes eligibility instead. Both respond to electricity conditions, but their service costs differ: Slurm waiting and DNN training-time penalties must be interpreted for their respective workloads.

\subsection{Estimators, measurements and validation}
Operational carbon combines energy with an electricity-emissions factor. Green Algorithms estimates energy from runtime, processor characteristics and allocated memory, assuming utilisation and facility overhead \citep{lannelongue2021green}. It works without direct measurements, but depends on representative coefficients. Carbontracker samples supported hardware interfaces and combines estimated energy with carbon information \citep{anthony2020carbontracker}. Its boundary matters: GPU monitoring does not automatically include the host CPU, storage, networking or cooling.

Ichnos combines scientific-workflow traces, resource power models and time-varying intensity to estimate carbon after execution \citep{west2024ichnos}. Its validation does not transfer to Stanage without comparable measurements. F-DATA combines Fugaku job, power and energy records, but cannot replace Stanage telemetry \citep{antici2025fdata}. Here, allocated resources are observed scheduling quantities used as activity proxies, not electrical-utilisation measurements.

A coefficient model permits unit and sensitivity checks; a meter-based estimator also permits comparison with independent electrical measurements. Trace estimators need credible power mappings and precise execution intervals. O'Brien and colleagues explain why predictive models depend on hardware and application characteristics~\citep{obrien2017survey}. Matching resource counts therefore establishes workload consistency, not the accuracy of estimated kilowatt-hours.

\subsection{Choosing a method with limited telemetry}
Green Algorithms needs workload descriptors; Carbontracker needs supported device measurements during execution~\citep{lannelongue2021green,anthony2020carbontracker}. Stanage provides descriptors and intervals but lacks usable per-job energy, making an explicit power scenario feasible retrospectively. Comparison with a future metered estimate would require matching the measurement scope, sampling interval and idle-energy allocation.

\subsection{Shared energy and the accounting boundary}
Kamiya and Coroam{\u a} explain how boundaries, source data and assumptions change data-centre energy estimates~\citep{kamiya2025datacentre}. A changed boundary can otherwise be mistaken for improved efficiency. Here, supplied power figures cover cage equipment, and allocation representations share that boundary. The review's global estimates provide context, not calibration for Stanage coefficients.

\subsection{Operational and lifecycle emissions}
ACT and the HPC lifecycle study by Li and colleagues separate operational emissions from those embodied in manufacture \citep{gupta2022act,li2023sustainable}. Moving work on existing hardware does not establish a smaller manufacturing footprint; this evaluation therefore covers operational electricity. Comparisons remain conditional on the power model. Shared assumptions aid comparison, but errors may not cancel when policies change resource placement.

The BLOOM study separates dynamic consumption, idle infrastructure and embodied emissions, each requiring different evidence and allocation assumptions~\citep{luccioni2023bloom}. Scheduling may change execution timing without changing manufacture or idle infrastructure. This dissertation distinguishes load-dependent and whole-cage carbon without transferring BLOOM's component ratios or training-utilisation assumptions to Stanage.

\section{Temporal flexibility and carbon information}

\subsection{Carbon exposure across a job's runtime}
Temporal shifting uses the fact that electricity carbon intensity varies during a job's feasible execution period. Wiesner and colleagues~\citep{wiesner2021wait} examine how flexibility windows, interruption assumptions and forecast accuracy affect this opportunity. For a non-preemptive job, the relevant exposure extends across its runtime. A favourable intensity at release can be followed by several unfavourable hours, so selecting the lowest single timestamp is only an approximation to the interval objective. This motivates comparing a simple time-shifting baseline with a policy that evaluates the complete candidate runtime.

\subsection{Flexibility, deadlines and queue constraints}
The workload limits this opportunity. Sukprasert and colleagues~\citep{sukprasert2024limitations} analyse ideal and practical shifting bounds across many regions, finding that attainable benefits depend on job duration, deadlines and service constraints. This supports keeping simple baselines in the comparison: a more complex rule should earn that complexity through a better trade-off. A large saving in one setting also need not generalise. In a university cluster, lower-carbon periods must coincide with available capacity and jobs whose users can tolerate a delay.

A flexibility window is different from an application deadline. A maximum admission delay limits the policy's own action, but cannot limit subsequent queueing or guarantee a completion time. Wiesner and colleagues explicitly consider the conditions under which work may be deferred or interrupted~\citep{wiesner2021wait}. The policies here use non-preemptive execution and a bounded change in release time. Since their control ends at admission, responsiveness must be measured after replay even when every decision stays within its budget.

An opportunity analysis asks whether the same work would have lower carbon exposure in a later interval. A scheduler evaluation goes further: it asks whether the proposed admission changes still reduce exposure once jobs compete for capacity. Sukprasert and colleagues show how practical restrictions can reduce ideal shifting potential~\citep{sukprasert2024limitations}. This project records the first comparison in the decision audit and the second in the completed schedule. The two records help explain small or negative realised changes without assuming that the calculation is wrong.

\subsection{Regional intensity and forecast availability}
NESO provides the regional signal used here. Its documentation describes electricity-generation CO$_2$ intensity, regional modelling and half-hour updates; the API identifies South Yorkshire as region 5 \citep{nesoCarbonIntensity,nesoAPI}. A regional value is neither a site electricity meter nor a complete lifecycle greenhouse-gas inventory. An additional distinction concerns information timing. A historical series obtained after an event does not by itself establish which forecast was available when a scheduling decision would have been made. Consequently, retrospective carbon exposure, forecast-informed decisions and empirical forecast-error evaluation require separate descriptions. Synthetic perturbations can test sensitivity, but they cannot substitute for matched archived forecasts and later estimates.

\subsection{Attributional and marginal emissions}
Another distinction concerns the emissions question being asked. Dodge and colleagues distinguish marginal-emissions accounting, which addresses the change associated with an intervention, from attributional accounting based on average electricity intensity~\citep{dodge2022carbon}. The NESO series used here is a regional generation-intensity signal. Multiplying modelled consumption by that signal attributes an operational footprint to the simulated schedule; it does not measure the additional generator that responds to a particular admission decision. The term ``carbon reduction'' in the results therefore refers to a paired reduction under the stated attributional model. It should not be read as a validated causal change in grid emissions.

\section{Carbon-aware scheduling and Slurm admission}

GREEN combines carbon objectives with cluster scheduling using ML-job energy information, temporal flexibility and resource allocation, with mechanisms intended to protect completion time \citep{xu2025green}. It uses active scheduling rounds, job scaling and checkpoint-based preemption. Stanage's mixed historical workload lacks equivalent energy profiles and evidence that arbitrary jobs can resize or restart safely. GREEN therefore informs the choice to consider full-execution carbon, resource pressure and waiting protection, but its reported savings are not a benchmark that admission-only control on Stanage should be expected to reproduce.

PCAPS makes another constraint explicit: delaying one task may block dependent tasks and extend the completion of an entire data-processing job. It combines precedence information with a configurable carbon-versus-completion-time objective \citep{lechowicz2025pcaps}. This supports transparent trade-off parameters, although the DAG information used by PCAPS is different from the eligibility and resource fields available here. The Stanage study preserves observed eligibility as the boundary after recorded dependencies or holds and does not claim to reconstruct the full application dependency graph.

CATS offers a lighter integration approach: a Python tool that chooses lower-carbon periods for task execution \citep{cats2025}. Its current documentation includes use with \texttt{sbatch} \citep{catsSchedulers}. Slurm's \texttt{--begin} option delays allocation eligibility, but does not promise a start at the specified time \citep{schedmdSbatch}. The same separation applies here. An external admission layer chooses the release, while Slurm determines the actual start and placement. Carbon must therefore be calculated from the realised simulated schedule, not inferred from the selected release alone.

Deeper changes to Slurm are also possible. Chadha and colleagues~\citep{chadha2020slurm} extend it to schedule malleable applications that can change their resource allocation at runtime. That work demonstrates a different intervention with different application requirements. Keeping the historical Stanage jobs rigid allows the present experiment to isolate admission timing, while leaving resizing, preemption and hardware power control outside the implemented policy family.

These schedulers control different variables. GREEN changes the amount and timing of ML resources, PCAPS considers precedence-constrained tasks, and CATS provides a lightweight task-timing interface~\citep{xu2025green,lechowicz2025pcaps,cats2025}. Their mechanisms are useful references, but their published percentages do not form a common benchmark: jobs, resource models and service constraints differ. Reproducing a reported percentage would require reproducing those conditions. This study tests selected ideas through a shared admission interface so that every policy controls the same part of the scheduling process.

\section{Waiting budgets and the distribution of service costs}

Carbon reduction becomes a service-quality question when it is obtained through delay. Starburst, although designed for hybrid-cloud cost control, offers a useful waiting-budget framework for exposing this trade-off \citep{luo2024starburst}. It assigns different waits to jobs with different resource demands and allows administrators to choose a cost-completion-time balance. The principle can inform a carbon-aware policy without importing its cloud-bursting assumptions. In particular, a bounded allowance can prevent a short job from receiving a delay that is disproportionately large relative to its execution time.

The Adaptive policy uses this idea by adding a runtime-related waiting budget, a minimum expected carbon benefit and an explicit waiting penalty to carbon-sensitive admission. These are project design choices informed by the literature, not a reproduction of Starburst or GREEN. The shared Slurm replay tests whether they improve the trade-off. Even when the admission layer respects its budget, reordering the workload can change the waiting of other jobs.

Aggregate waiting measures can hide important differences. The mean may fall while the tail grows, and an apparently favourable percentile may conceal delay concentrated on one user. The evaluation therefore reports waiting percentiles, slowdown and user-level delay alongside modelled carbon. These measures show the balance among carbon, responsiveness and fairness. A setting can be useful for one objective and costly for another, which a single success or failure label would miss.

Waiting budgets do not make unlike objectives interchangeable. Starburst's cost-versus-completion-time control motivates explicit service preferences, while PCAPS illustrates why the completion effects of delaying a task depend on surrounding work~\citep{luo2024starburst,lechowicz2025pcaps}. A weighted policy score is consequently a decision rule, not a complete evaluation metric. Carbon, total waiting and user burden must still be reported in their original units. This permits a reader to see whether a favourable score was obtained through a material service cost and whether different users experienced that cost unevenly.

User-level carbon attribution introduces a further, separate choice. Wassermann and colleagues distinguish operator-level reporting from user-centred accounting and make the allocation of idle-resource carbon explicit~\citep{wassermann2024usercentric}. Their framework connects measurement and resource use to accounting responsibilities. By contrast, the present user metrics allocate deliberate waiting, not a carbon bill. A user who receives much of the policy delay is not necessarily responsible for a corresponding share of cage emissions. Maintaining that distinction allows burden concentration to be discussed without assigning an unmeasured individual environmental footprint.

\section{Trace-driven simulation and experimental credibility}

The Slurm simulator lineage offers a way to examine scheduling changes without intervening in production. The earlier simulator retains a real Slurm instance with modifications for historical-job replay and faster execution \citep{simakov2018simulator}. The later work by Simakov and colleagues~\citep{simakov2022accurate} emphasises careful transformation of Slurm code and the variability of scheduling realisations. Separately, Jokanovic and colleagues~\citep{jokanovic2018validation} compare a Slurm simulator with a physical installation. Such studies support the evaluation approach while also showing why fidelity must be established for a particular simulator, configuration and workload.

The scheduling configuration can affect outcomes even when workload membership is fixed. Simakov and colleagues study how node sharing and controller configuration change scheduler responsiveness and waiting in Slurm replay~\citep{simakov2018sharing}. This makes placement and sharing relevant to both service and occupancy analysis. A sum of requested nodes, for example, does not establish the number of distinct hosts active when jobs can share nodes. The current sensitivity calculations retain these different views, while the experiment keeps its chosen Slurm configuration fixed across policies.

Slurm backfill considers requested runtime and resources when deciding whether lower-priority work can proceed without delaying higher-priority reservations \citep{schedmdScheduling}. Admission changes can therefore alter interactions among jobs, even when each job's resources and runtime remain fixed. The current Stanage documentation confirms a heterogeneous CPU/GPU environment and Slurm-based submission, but it describes the current service and cannot alone establish the complete 2025 hardware inventory or production configuration \citep{sheffieldStanage}. Site history, reservations and administrator actions remain separate evidence requirements.

The comparison here controls job membership, resource requests, runtimes, reconstructed carry-in state and the carbon-accounting horizon, and requires complete execution. These controls make policy outcomes easier to interpret, but cannot establish that simulated waiting matches production waiting. A structurally valid run may also show no carbon benefit. Repeated baselines and checks across dates are especially important when estimated policy differences are small.

Simulator validity has several layers. Input checks establish that the intended job population and resource requests were encoded correctly. Execution checks establish that those jobs were submitted, scheduled and completed. Comparison with a physical system addresses a further question: whether the simulator reproduces the scheduling behaviour of that system under comparable conditions. Jokanovic and colleagues conduct this latter kind of validation for their simulator setting~\citep{jokanovic2018validation}. Passing the first two layers here does not inherit the third from their paper.

The simulator version and configuration also affect the experiment. The earlier Slurm simulator study examines implementation and parameter effects, while the later work identifies variation between scheduling realisations~\citep{simakov2018simulator,simakov2022accurate}. Accordingly, this project saves configuration and workload hashes with its results and repeats each date's baseline. A small carbon difference that changes sign between repeats is weaker evidence than a difference outside the observed repeat range. Chapter~\ref{ch:evaluation} makes this comparison using the saved outcomes.

\section{Planning and reporting a computational study}

The evaluation process needs a clear record of what was planned and what was later explored. Souter and colleagues recommend planning analyses to reduce avoidable reruns while retaining useful replication~\citep{souter2023recommendations}. Here, pilot windows informed policy development before the multi-date campaign was frozen. Recording that plan defines the intended comparison, but does not turn earlier development data into an untouched test set. Chapter~\ref{ch:evaluation} therefore separates historical diagnostics, the completed main comparison and the later bounded engineering extension.

Transparent reporting connects a footprint estimate to its calculation method and execution context~\citep{souter2023recommendations}. For this study, that requires identifying the source workload, policy parameters, power scenario, carbon series and accounting horizon. A second accounting problem would be the footprint of running the research experiments themselves. The simulated Stanage footprint is a counterfactual output and cannot substitute for the actual electricity consumed by the machine hosting the simulator. No estimate of that separate experimental footprint is claimed without the corresponding execution and energy records.

\section{Synthesis and research position}

These studies motivate explicit energy boundaries, full-runtime exposure, capacity awareness, waiting safeguards and controlled replay. This dissertation tests these principles together on a mixed university workload.

Table~\ref{tab:related-work-boundary} brings these distinctions together. The useful comparison is the intervention each system can make and the information it requires. This makes the scope of an admission-only experiment visible before its results are interpreted.
\begin{table}[htbp]
\centering
\small
\caption[Related work and its relevance to admission control]{Selected related work compared by intervention and information needs. The final column states the design relevance to this project; it is not a comparison of measured savings across systems.}
\label{tab:related-work-boundary}
\setlength{\tabcolsep}{4pt}
\renewcommand{\arraystretch}{1.2}
\begin{tabular}{>{\raggedright\arraybackslash}p{0.18\linewidth}>{\raggedright\arraybackslash}p{0.36\linewidth}>{\raggedright\arraybackslash}p{0.38\linewidth}}
\toprule
Work & Intervention and information & Relevance to the Stanage study \\
\midrule
Let's Wait Awhile~\citep{wiesner2021wait} & Execution timing within a flexibility window; job duration and carbon information. & Evaluate exposure over runtime and state the delay allowance and forecast assumptions. \\
GREEN~\citep{xu2025green} & ML-cluster scheduling with energy profiles, scaling and checkpoint-based preemption. & Consider carbon and capacity together; application resizing and checkpointing are outside the implemented admission layer. \\
PCAPS~\citep{lechowicz2025pcaps} & Carbon-aware scheduling with task-precedence information and a completion-time trade-off. & Make service costs explicit; observed eligibility is available here, but a complete task DAG is not. \\
Starburst~\citep{luo2024starburst} & Job-sensitive waiting budgets for hybrid-cloud cost and completion-time control. & Motivate bounded waiting; cloud migration and cost optimisation are not part of this experiment. \\
CATS~\citep{cats2025,catsSchedulers} & Lightweight carbon-aware task timing, including documented \texttt{sbatch} integration. & Separate admission eligibility from the actual allocation and start chosen by Slurm. \\
This study & Fixed-job release-time policies with declared carbon/load information, carry-in and matched Slurm replay. & Compare modelled carbon, waiting and user burden. Multi-day outcomes remain to be completed. \\
\bottomrule
\end{tabular}
\end{table}


This dissertation compares normal admission, B1 shifting, B2 gating and Adaptive admission with explicit waiting protection. Oracle and history-based variants add comparisons under different information and safeguards. The completed multi-day evaluation examines how consistently modelled carbon, waiting and user burden change together. Chapter~\ref{ch:evaluation} presents the results, including negative and no-action cases.
